% !TEX program = xelatex

\documentclass{article}
\usepackage[UTF8]{ctex}
\usepackage[tc]{titlepic}
\usepackage{titlesec}
\usepackage{cite}
\usepackage{fancyhdr}
\usepackage{booktabs}
\usepackage{graphicx}
\usepackage{geometry}
\usepackage[section]{placeins}
\geometry{a4paper,scale=0.8}
\pagestyle{fancy}

\lhead{作业 2——矩阵分解\\\today}
\chead{中国科学技术大学\\数学建模课程}
\rhead{SVD 图像压缩}

\title{基于奇异值分解的图像压缩}
\author{姓名：胡广理 \\ 学号：PB24010500}

\begin{document}
\maketitle
\begin{abstract}
    本文实现了基于奇异值分解（SVD）的灰度图像压缩。遵循课程要求，SVD 的计算不直接调用高阶 SVD 接口，而是通过对称矩阵的特征分解实现。对仓库示例图片进行秩-$k$ 重构，使用 PSNR 和（全局）SSIM 评估压缩效果，并讨论方法的复杂度与改进方向。
\end{abstract}


\setcounter{section}{1}
\section*{一、前言（问题的提出）}
图像压缩是减小存储与传输代价的常见任务。SVD 为矩阵提供低秩近似，常用于图像近似与压缩：保留重要奇异值即可恢复主要视觉信息。

\setcounter{section}{3}
\section*{三、问题分析}
给定灰度图像矩阵 $A\in \mathbb{R}^{m\times n}$，目标是构造秩-$k$ 近似 $A_k$，使得在保持视觉质量的同时显著压缩数据量。我们将比较不同 $k$ 值的重构质量与代价。

\setcounter{section}{6}
\section*{六、数学模型建立}
令 $A=U\Sigma V^T$ 为 $A$ 的奇异值分解，其中 $\Sigma=\mathrm{diag}(\sigma_1,\dots,\sigma_r)$，$r=\min(m,n)$，奇异值按非增顺序排列。秩-$k$ 重构为
$$A_k=\sum_{i=1}^k\sigma_i u_i v_i^T = U_k \Sigma_k V_k^T.$$ 

在计算上，当 $m\ge n$，我们先计算 $A^TA$ 的特征分解 $A^TA=V\Lambda V^T$（对称矩阵，可用 \texttt{eigh} 稳定计算），取 $\sigma_i=\sqrt{\lambda_i}$，再计算 $u_i=(A v_i)/\sigma_i$。若 $n>m$，对 $AA^T$ 做类似处理以得到 $U$，再求 $V$。该方法避免直接调用 \texttt{linalg.svd}，但依然依赖 \texttt{linalg.eigh} 求对称特征分解。

\subsection*{时间复杂度与存储}
完整特征分解代价约 $O(\min(m,n)^3)$，适用于中小尺寸图像。存储方面，秩-$k$ 近似只需保存 $U_k( m\times k)$、$\Sigma_k(k)$ 与 $V_k^T(k\times n)$，总体元素数为 $k(m+n+1)$。

\setcounter{section}{7}
\section*{七、结果（与对比）}
\subsection*{实验设置}
实验使用仓库示例图片 \texttt{imgs/CMakeTools.png}（转为灰度）作为输入，在不同 $k$ 值下重构并评估 PSNR 与全局 SSIM（实现见代码）。实验在 \texttt{hw\_2/option1} 目录下运行，重构图与汇总图保存在 \texttt{demo\_out}。

\subsection*{可视化结果}
下面展示汇总对比图与各 $k$ 的重构结果。
\begin{figure}[h]
    \centering
    \includegraphics[width=0.85\textwidth]{demo\_out/summary.png}
    \caption{原图与秩-$k$ 重构汇总（k=5,20,50）}
\end{figure}

\begin{figure}[h]
    \centering
    \includegraphics[width=0.32\textwidth]{demo\_out/recon\_k5.png}
    \includegraphics[width=0.32\textwidth]{demo\_out/recon\_k20.png}
    \includegraphics[width=0.32\textwidth]{demo\_out/recon\_k50.png}
    \caption{左：k=5；中：k=20；右：k=50}
\end{figure}

\subsection*{定量指标}
\begin{table}[h]
\centering
\caption{不同 $k$ 值下的 PSNR 与 SSIM}
\begin{tabular}{ccc}
\toprule
k & PSNR (dB) & SSIM \\
\midrule
5 & 20.417 & 0.8197 \\
20 & 24.997 & 0.9431 \\
50 & 32.372 & 0.9901 \\
\bottomrule
\end{tabular}
\end{table}

\setcounter{section}{8}
\section*{八、结论}
实验表明：随着保留秩 $k$ 的增加，重构图像的视觉质量与 PSNR/SSIM 均显著提升。对该示例而言，$k=20$ 已能达到较好的视觉保真度，而 $k=50$ 接近原图。

\setcounter{section}{9}
\section*{九、问题与扩展}
当前实现受限于完整特征分解的计算开销，后续可考虑：
\begin{itemize}
    \item 使用截断或随机化 SVD（randomized SVD）以在大图上加速；
    \item 对彩色图像按通道处理或使用分块 SVD（block-wise SVD）；
    \item 使用滑动窗口实现标准 SSIM 以获得更精确的感知质量度量；
    \item 提供交互式 GUI 便于可视化比较与参数调优。
\end{itemize}

\bibliographystyle{ieeetr}
\bibliography{refer}

\end{document}
